% The mathematics of LinearDecisionRules.jl 0.4.
%
% Build: `latexmk` in this folder (see latexmkrc); the PDF is written to build/.
% Source conventions: semantic line breaks (one sentence, or one clause, per line),
% notation through the macros of macros.tex.
\documentclass[11pt,a4paper]{article}

\usepackage[T1]{fontenc}
\usepackage[utf8]{inputenc}
\usepackage{lmodern}
\usepackage[margin=1in]{geometry}
\usepackage{microtype}
\usepackage{amsmath,amssymb,amsthm,mathtools}
\usepackage{booktabs}
\usepackage{enumitem}
\usepackage[numbers,sort&compress]{natbib}
\usepackage{xcolor}
\usepackage{url}
\usepackage[colorlinks=true,linkcolor=blue!50!black,citecolor=green!40!black,urlcolor=blue!50!black]{hyperref}
\usepackage[capitalise,noabbrev]{cleveref}
\usepackage{thmtools}

% Notation of the document.
% One macro per recurring object, so that a change of notation is a change of one line.

% Sets and operators
\newcommand{\R}{\mathbb{R}}
\newcommand{\N}{\mathbb{N}}
\newcommand{\E}{\mathbb{E}}
\newcommand{\Prob}{\mathbb{P}}
\newcommand{\T}{\top}
\newcommand{\tr}{\operatorname{tr}}
\newcommand{\Cov}{\operatorname{Cov}}
\newcommand{\conv}{\operatorname{conv}}
\newcommand{\supp}{\operatorname{supp}}
\newcommand{\blockdiag}{\operatorname{blockdiag}}
\newcommand{\inner}[2]{\langle #1, #2 \rangle}
\newcommand{\norm}[1]{\lVert #1 \rVert}
\newcommand{\abs}[1]{\lvert #1 \rvert}
\newcommand{\one}{\mathbf{1}}
\newcommand{\ind}[1]{\mathbb{1}\{#1\}}
\newcommand{\eye}{I}
\newcommand{\eone}{e_1}
\newcommand{\opt}{\operatorname{opt}}
\newcommand{\val}{\operatorname{val}}
\newcommand{\defeq}{\coloneqq}

% Uncertainty and its law
\newcommand{\xis}{\xi}                 % ξ = (1, η)
\newcommand{\etas}{\eta}               % the uncertainties
\newcommand{\Xis}{\Xi}                 % the support set
\newcommand{\Xib}[1]{\Xi_{#1}}         % one support block
\newcommand{\law}{\Prob}               % the law of η
\newcommand{\mub}{\bar\mu}             % E[ξ] = (1, μ)
\newcommand{\Mom}{M}                   % E[ξ ξᵀ]
\newcommand{\Cv}{\Sigma}               % covariance of ξ, block diagonal
\newcommand{\Cvb}[1]{\Sigma_{#1}}      % one covariance block
\newcommand{\Wp}{W'}                   % rows of Ξ as a cone, W'ξ ≥ 0
\newcommand{\nvis}{v}                  % number of observed columns
\newcommand{\Jr}{J}                    % structural columns of a row
\newcommand{\Touch}{\mathcal{B}}       % support blocks touched by a row
\newcommand{\Rows}{\mathcal{R}}        % the inequalities
\newcommand{\Rules}{\mathcal{X}}       % the rules allowed by the information structure
\newcommand{\stage}{\tau}              % stage of a column of ξ
\newcommand{\dstage}{t}                % stage of a decision

% Liftings
\newcommand{\Lift}{L}
\newcommand{\Retr}{R}
\newcommand{\lxi}{\tilde\xi}
\newcommand{\leta}{\tilde\eta}
\newcommand{\lXi}{\tilde\Xi}
\newcommand{\lX}{\tilde X}
\newcommand{\lM}{\tilde M}
\newcommand{\lk}{\tilde k}
\newcommand{\lm}{\tilde m}

% Values of the problems
\newcommand{\Vstar}{V^{\star}}
\newcommand{\Vp}{V^{\mathrm{P}}}
\newcommand{\Vd}{V^{\mathrm{D}}}
\newcommand{\Vs}{V^{\mathrm{S}}_N}

% Problems
\newcommand{\Pb}{\textup{(P)}}
\newcommand{\PbLDR}{\textup{(P-LDR)}}
\newcommand{\DbLDR}{\textup{(D-LDR)}}
\newcommand{\SbLDR}{\textup{(S-LDR)}}

% Chi-squared distribution function and quantile
\newcommand{\Fchi}[1]{F_{\chi^2_{#1}}}

% Software
\newcommand{\pkg}{\textsf{LinearDecisionRules.jl}}
\newcommand{\code}[1]{\texttt{#1}}
\newcommand{\file}{\path} % reads its argument verbatim, so `_` needs no escape


% Statements (declared with thmtools, so that cleveref names the environments that share a counter)
\declaretheorem[name=Theorem,numberwithin=section]{theorem}
\declaretheorem[name=Proposition,sibling=theorem]{proposition}
\declaretheorem[name=Lemma,sibling=theorem]{lemma}
\declaretheorem[name=Corollary,sibling=theorem]{corollary}
\declaretheorem[name=Definition,sibling=theorem,style=definition]{definition}
\declaretheorem[name=Assumption,sibling=theorem,style=definition]{assumption}
\declaretheorem[name=Remark,sibling=theorem,style=remark]{remark}
\declaretheorem[name=Example,sibling=theorem,style=remark]{example}

\numberwithin{equation}{section}

\title{The mathematics of \pkg{} 0.4\\[0.5ex]
\large Primal, dual and sampled linear decision rules, liftings,\\
and how the reformulated problems are kept small}
\author{Bernardo Freitas Paulo da Costa\thanks{Escola de Matem\'atica Aplicada, Funda\c{c}\~ao Getulio Vargas (FGV), Rio de Janeiro, Brazil.}
\and Joaquim Dias Garcia\thanks{PSR, Rio de Janeiro, Brazil \& Soma Energy.}}
\date{Draft, October 2026}

\begin{document}

\maketitle

\begin{abstract}
\pkg{} reformulates a multistage stochastic program written in JuMP, with uncertain right-hand sides and costs, into deterministic problems over the coefficients of linear decision rules.
The primal approximation restricts the policies to affine functions of the observed uncertainties and gives a feasible policy, hence a bound on the optimal value from the side of the objective;
the dual approximation relaxes the constraints to hold in expectation against the functions that describe the support, and gives a bound from the other side.
Liftings extend both to piecewise affine rules, and a sampled formulation enforces the constraints on drawn outcomes only.
These notes describe the mathematics as implemented in version 0.4,
and the computational choices that keep the reformulated problems small:
the information structure as prefixes of the sorted uncertainties,
the support as a product of blocks,
the second-moment matrix as a rank-one term plus a block-diagonal covariance,
and sparse retractions for the liftings.
We also describe how laws are conditioned and their moments estimated,
and the confidence-ellipsoid normal law that is being ported from the previous version.
\end{abstract}

\tableofcontents

\section{Introduction}
\label{sec:introduction}

To be written \citep{KWG2011,GWK2015}.

\section{The problem and its notation}
\label{sec:problem}

To be written.

\section{Linear decision rules}
\label{sec:rules}

\subsection{The restricted problem}

A linear decision rule is $x(\xi) = X\xi$ for a matrix $X \in \R^{n\times k}$.
It respects the information structure if row $i$ of $X$ vanishes beyond the observed prefix:
\[
\Rules \defeq \{\, X \in \R^{n\times k} : X_{ij} = 0 \text{ for } j > \nvis_i \,\}.
\]
Nonanticipativity is thus a sparsity pattern, and the package only creates the variables $X_{ij}$, $j \le \nvis_i$ (\code{\_rule\_variables!}, \file{src/formulations/formulation.jl}):
$\sum_i \nvis_i$ variables, instead of $nk$.
For an integer decision, $X_{i1}$ is an integer variable.
Restricting \eqref{eq:P} to these rules, and asking the constraints to hold on $\Xis$, gives the semi-infinite problem
\begin{equation}
\label{eq:PLDR}
\tag{P-LDR}
\begin{aligned}
\opt_{X \in \Rules} \quad & \E\bigl[ \xi^\T X^\T P X \xi + \xi^\T X^\T C \xi \bigr] + \E\bigl[\eta^\T Q\eta + d^\T\eta + f\bigr] \\
\text{s.t.} \quad & (A_e X - B_e)\xi = 0, \quad (B_u - A_u X)\xi \ge 0, \quad (A_l X - B_l)\xi \ge 0, \\
& \ell \le X\xi \le u \qquad \forall \xi \in \Xis.
\end{aligned}
\end{equation}
Its equalities and its objective need no approximation;
we treat them here, and the inequalities in the next three sections.

\subsection{Equalities: coefficient matching}

\begin{proposition}[Coefficient matching]
\label{prop:matching}
Suppose that $\Xis$ has a nonempty interior in the hyperplane $\{\xi_1 = 1\}$.
Then $(A_e X - B_e)\xi = 0$ for all $\xi \in \Xis$ if and only if $A_e X = B_e$.
For a row $r$ of $A_e$ and $X \in \Rules$, this amounts to the $\abs{\Jr_r}$ scalar equalities $(a_r^\T X)_j = b_{rj}$, $j \in \Jr_r$.
\end{proposition}

\begin{proof}
A linear function that vanishes on a set vanishes on its linear span, which is $\R^k$ when $\Xis$ has a nonempty relative interior in $\{\xi_1 = 1\}$.
For $j \notin \Jr_r$ both sides of $(a_r^\T X)_j = b_{rj}$ are zero by definition of $\Jr_r$.
\end{proof}

This is how \citet[Section~2.2]{KWG2011} eliminate the equalities $AX + S = B$.
The package adds them once per structural column (\code{\_add\_equalities!}), in all three formulations.
If a row has $b_{rj} \neq 0$ for some $j > \nvis_r$,
no rule of its decisions can match the coefficient of the unobserved column $j$, and the model is rejected when it is read:
this is the case of a balance equation that involves an uncertainty that none of its decisions observes, typically a modelling error in the stages.

\subsection{The objective through the structured moments}

With $x = X\xi$, the expected objective is a quadratic function of $X$ that depends on the law only through $\mub$ and $\Cv$.

\begin{proposition}[Objective]
\label{prop:objective}
For $X \in \Rules$, and with $X_\ell$, $C_\ell$ the columns of $X$ and $C$ in the covariance block $\ell$,
\begin{align}
\E[x^\T P x] &= (X\mub)^\T P (X\mub) + \sum_\ell \inner{X_\ell^\T P X_\ell}{\Cvb{\ell}}, \label{eq:obj-quadratic}\\
\E[x^\T C \xi] &= (C\mub)^\T (X\mub) + \sum_\ell \inner{C_\ell \Cvb{\ell}}{X_\ell}, \label{eq:obj-linear}\\
\E[\eta^\T Q\eta + d^\T\eta + f] &= f + d^\T\mu + \mu^\T Q \mu + \sum_\ell \inner{Q_{\ell\ell}}{\Cvb{\ell}}, \label{eq:obj-constant}
\end{align}
where $\inner{A}{B} = \tr(A^\T B)$ and $Q_{\ell\ell}$ is the diagonal block of $Q$ on the coordinates of block $\ell$.
\end{proposition}

\begin{proof}
$\E[x^\T P x] = \tr(X^\T P X \Mom)$ and $\E[x^\T C\xi] = \tr(C^\T X \Mom) = \inner{C}{X\Mom}$;
substitute \eqref{eq:structured-moments}, use that $\Cv$ is block diagonal and that $\Cvb{\ell}$ is symmetric.
The third line is the same computation with $\eta$.
\end{proof}

The package never forms $\Mom$ nor a product of dense matrices:
$C\mub$ is one sparse product, $C_\ell\Cvb{\ell}$ involves only the rows of $C$ with a nonzero in block $\ell$,
and the products $X_\ell$ are restricted to the columns that each decision observes,
visited in increasing order so that the loop over a block stops at the first unobserved column (\code{\_set\_objective!}).
In particular, a linear objective ($P = 0$) is a linear function of $X$ with at most $\sum_i \nvis_i$ terms.
When the objective is quadratic,
the problem \eqref{eq:PLDR} is convex if $P \succeq 0$ for a minimization ($P \preceq 0$ for a maximization),
and this is also what the dual bound of \cref{sec:dual} needs;
the package passes the quadratic objective to the inner solver without checking it.

\subsection{Inequalities and bounds}

All inequalities of \eqref{eq:PLDR} have the same form.
For a row $r$ of $(A_u, B_u)$ let $\sigma_r = 1$, for a row of $(A_l, B_l)$ let $\sigma_r = -1$, and write
\begin{equation}
\label{eq:slack-rule}
s_r(\xi) \defeq \sigma_r \bigl(b_r^\T - a_r^\T X\bigr)\xi \eqqcolon S_r \xi \ge 0 \qquad \forall \xi \in \Xis.
\end{equation}
The row vector $S_r$ is the slack rule of \citet{KWG2011}; by \eqref{eq:structural-columns}, $S_{rj} = 0$ for $j \notin \Jr_r$, whatever $X \in \Rules$.
The bounds of the decisions are inequalities of the same form, with $a_r = e_i$ and $b_r = \ell_i e_1$ or $u_i e_1$,
except for the decisions that observe only the constant ($\nvis_i = 1$):
their rule is a number $X_{i1}$, and their bounds are simple bounds on it.
We denote by $\Rows$ the set of all these inequalities (\code{\_inequalities}).

A deterministic inequality ($\Jr_r = \{1\}$) is the plain linear inequality $S_{r1} \ge 0$ in all formulations.
The others are semi-infinite, and the three formulations differ in how they treat them:
the primal approximation enforces them on all of $\Xis$ (\cref{sec:primal}),
the dual approximation relaxes them to hold in expectation against the functions that describe $\Xis$ (\cref{sec:dual}),
and the sampled formulation enforces them on finitely many outcomes (\cref{sec:sampled}).

\section{The primal approximation}
\label{sec:primal}

To be written.

\section{The dual approximation}
\label{sec:dual}

\subsection{The relaxation and why it bounds}

\citet[Section~2.3]{KWG2011} obtain a lower bound on a minimization by dualizing the equalities $Ax + s = b$ of the problem with slacks,
restricting the dual multipliers to linear rules $y(\xi) = Y\xi$, and carrying out the maximization over $Y$.
What remains is the original problem with the equalities required to hold only in expectation against $\xi$,
$\E[(Ax(\xi) + s(\xi) - b(\xi))\xi^\T] = 0$, and $s(\xi) \ge 0$ almost surely.
Its tractable form \citep[(2.8) and (4.6)]{KWG2011} replaces the policies by their projections $X$, $S$ on the linear functions of $\xi$,
defined by $X\Mom = \E[x(\xi)\xi^\T]$ and $S\Mom = \E[s(\xi)\xi^\T]$,
and the almost sure nonnegativity of $s$ by a finite set of linear inequalities on $S\Mom$, one per row of $W$.

We write the test functions as a cone:
let $\Wp$ be the matrix with the rows $\eone^\T$ and $(-h_{bp}, W_{bp})$ for every row $p$ of every support block $b$, embedded in the $k$ columns of $\xi$,
so that $\Wp\xi = (\xi_1,\ W\eta - h\xi_1)$ and
\[
\xi \in \Xis \implies \Wp\xi \ge 0.
\]
In \citet[(2.1b)]{KWG2011}, $W = (\eone, -\eone, \hat W^\T)^\T$ and $h = (1, -1, 0, \dots, 0)$, so that $W - h\eone^\T$ consists of two zero rows and the rows $\hat W$;
the rows of $\Wp$ other than $\eone^\T$ play the role of $\hat W$, with the constants of $h$ moved to the column of $\xi_1$.
The dual approximation solved by the package is
\begin{equation}
\label{eq:DLDR}
\tag{D-LDR}
\begin{aligned}
\opt_{X \in \Rules} \quad & \text{objective of \cref{prop:objective}} \\
\text{s.t.} \quad & A_e X = B_e, \\
& S_r = \sigma_r(b_r^\T - a_r^\T X) \\
& S_{r1} \ge 0 \text{ for the deterministic } r \in \Rows,\\
& S_r \Mom \Wp^\T \ge 0 \quad \text{for the other } r \in \Rows.
\end{aligned}
\end{equation}
The inequalities $S_r\Mom\Wp^\T \ge 0$ are $\E[s_r(\xi)\,(\Wp\xi)^\T] \ge 0$ for linear $s_r$:
the slack must be nonnegative ``on average against every function that is nonnegative on $\Xis$ by construction''.

\begin{theorem}[Dual bound]
\label{thm:dual-bound}
Suppose \cref{ass:support} and the following conditions hold.
\begin{enumerate}[label=(\roman*)]
\item\label{it:lce} \emph{Linear conditional expectations}: for every stage $t$ there is a matrix $N_t$ such that
$\E[\xi \mid \xi_{1:\nvis(t)}] = N_t\, \xi_{1:\nvis(t)}$ almost surely.
\item\label{it:convex} $P \succeq 0$ if $\opt = \min$, and $P \preceq 0$ if $\opt = \max$.
\end{enumerate}
Then for every feasible policy $x$ of \eqref{eq:SP} with square-integrable components, the matrix $X = \E[x(\xi)\xi^\T]\Mom^{-1}$ is feasible in \eqref{eq:DLDR},
and its objective is at least as good as that of $x$.
Consequently, for a minimization, $\val\eqref{eq:DLDR} \le \val\eqref{eq:SP} \le \val\eqref{eq:PLDR}$, and the reverse for a maximization.
\end{theorem}

\begin{proof}[Proof sketch]
The equalities hold almost surely, so in expectation against $\xi^\T$, which gives $A_eX\Mom = B_e\Mom$, hence $A_eX = B_e$.
For an inequality, $s_r(\xi) \ge 0$ and $\Wp\xi \ge 0$ almost surely, so $\E[s_r(\xi)(\Wp\xi)^\T] \ge 0$,
and $\E[s_r(\xi)\xi^\T] = S_r\Mom$ with $S_r = \sigma_r(b_r^\T - a_r^\T X)$ by linearity.
The objective is linear in $\E[x\xi^\T] = X\Mom$, except for the term in $P$:
writing $x = X\xi + \varepsilon$ with $\E[\varepsilon\xi^\T] = 0$, the cross term vanishes and $\E[\varepsilon^\T P \varepsilon]$ has the sign of $P$.
It remains to show that $X \in \Rules$, which is where condition~\ref{it:lce} enters:
the projection of a function of the history $\xi_{1:\nvis_i}$ on the linear functions of the whole of $\xi$ only involves the history if the conditional expectation of the future given the history is linear,
which is \citet[Lemma~2]{KWG2011}.
An integer decision is a constant, and its projection is the same constant, so integrality is preserved.
The complete argument is in \cref{app:dual}.
\end{proof}

This is the lower bound of \citet[Section~4]{KWG2011}, whose tractable form is their Theorem~3, under their condition $\E_t[\xi] = M_tP_t\xi$,
with a quadratic objective and integer first-stage decisions added.
They also show that \eqref{eq:DLDR} is exact for their relaxed problem when the original one is strictly feasible \citep[Proposition~4]{KWG2011};
the bound itself does not need strict feasibility.
\citet{GWK2015} assume instead that the stages are independent \citep[condition (M4)]{GWK2015},
which implies condition~\ref{it:lce} \citep[footnote~3]{KWG2011}.

% What condition (i) implies?  In this "no stage for constraints", is it needed?
\begin{remark}[What condition~\ref{it:lce} asks of the laws]
\label{rem:lce}
Distinct declarations are independent, so condition~\ref{it:lce} holds as soon as each declaration lies in one stage,
and more generally when the law of each declaration has linear conditional expectations across its own stages.
A multivariate normal law has this property, and so has its truncation to a confidence ellipsoid (\cref{sec:confidence-normal});
a discrete joint law of several stages generally does not.
The package checks condition~\ref{it:lce} law by law, for the laws whose components are in several stages (\code{\_linear\_conditional\_expectations}, \file{src/probability/moments.jl}):
it holds for a product of univariate laws (independent components), a normal law and \code{ConfidenceMvNormal};
for an \code{MvDiscreteNonParametric} it is checked exactly, for each stage $t$, by fitting an affine function of the history to the conditional means of the later components, one per distinct history, which must leave no residual;
for any other law, such as a law conditioned by rejection, it is unknown.
When it fails or is unknown for some law, the package warns, and marks the bound as not certified (\code{\_stage\_dependent\_laws}, \file{src/formulations/dual.jl}).
Correlation is not the test: a correlated law may satisfy condition~\ref{it:lce}, such as the truncated normal law above,
and an uncorrelated one may not, such as $\eta_2 = \eta_1^2$ with $\eta_1$ uniform on $\{-1, 0, 1\}$.
\end{remark}

\begin{remark}[The row $\eone^\T$]
\label{rem:e1-row}
In \citet{KWG2011}, $\Xis$ is compact, and the cone $\{z : (W - h\eone^\T)z \ge 0\}$ already implies $z_1 \ge 0$ \citep[proof of Proposition~3]{KWG2011};
the rows $\pm\eone^\T$ of their $W$ become zero rows of $W - h\eone^\T$.
The package allows unbounded blocks, and keeps the row $\eone^\T$ of $\Wp$ explicitly,
that is, the inequality $\E[s_r(\xi)] \ge 0$.
Without it, a model in which no uncertainty is bounded would have no inequality left in its dual approximation;
the previous version of the package had this defect.
\end{remark}

The relaxation depends on the set $\Xis$, not on the inequalities that describe it.

\begin{proposition}[Invariance under the description of $\Xis$]
\label{prop:description}
Adding to $W\eta \ge h$ an inequality that is valid on $\Xis$ does not change the feasible set of \eqref{eq:DLDR}.
In particular, redundant rows and positive scalings of rows have no effect on the dual bound.
\end{proposition}

\begin{proof}
Let $w^\T\eta \ge g$ be valid on the nonempty polyhedron $\{W\eta \ge h\}$.
By the affine form of Farkas' lemma there are $c \ge 0$ and $c_0 \ge 0$ with $w = W^\T c$ and $g = h^\T c - c_0$.
Then $w^\T\eta - g\xi_1 = c^\T(W\eta - h\xi_1) + c_0\xi_1$ is a nonnegative combination of the rows of $\Wp$,
and the new inequality $S_r\Mom(-g, w)^\T \ge 0$ is the same combination of inequalities of \eqref{eq:DLDR}.
\end{proof}

In the language of \citet[Proposition~3]{KWG2011}, the cone defined by the rows of $\Wp$ is the closed convex cone generated by $\Xis$.
The row $\eone^\T$ is needed for the proposition when $\Xis$ is unbounded.

\subsection{The structured form}

Written literally, $S_r\Mom\Wp^\T$ is a row with one entry per row of $\Wp$, each a sum of $k$ products of an entry of $S_r$, itself an expression of $X$, by an entry of the dense matrix $\Mom\Wp^\T$.
This is what made the dual of the previous version explode.
Substituting \eqref{eq:structured-moments},
\begin{equation}
\label{eq:dual-identity}
S_r \Mom \Wp^\T = (S_r\mub)(\Wp\mub)^\T + S_r\Cv\Wp^\T.
\end{equation}
Introduce the scalar $u_r \defeq S_r\mub$, the expected slack.
The row $\eone^\T$ of $\Wp$ gives $u_r \ge 0$, since $\Cv\eone = 0$ and $\eone^\T\mub = 1$.
The row $p$ of support block $b$ gives
\begin{equation}
\label{eq:dual-cone-row}
u_r\,(W_{bp}\mu_b - h_{bp}) + S_r\, \Cv_{:,b} W_{bp}^\T \ge 0,
\end{equation}
where $W_{bp}\mu_b - h_{bp} \ge 0$ because the mean lies in $\Xib{b}$,
and $\Cv_{:,b}$ are the columns of $\Cv$ of the coordinates of block $b$.
The matrix $V_b \defeq \Cv_{:,b}W_b^\T$ is computed once per support block (\code{\_dual\_cones}),
and its row $j$ is nonzero only if column $j$ shares a covariance block with some coordinate of $b$.
In \eqref{eq:dual-cone-row}, $S_r$ has nonzeros only in $\Jr_r$, so:
\begin{itemize}[nosep]
\item only the support blocks that share a covariance block with a structural column of $r$ give inequalities with a slack term;
\item a row of $\Wp$ whose term $S_r V_{b,:p}$ is structurally zero reduces to $u_r(W_{bp}\mu_b - h_{bp}) \ge 0$, implied by $u_r \ge 0$, and is dropped:
all the blocks that the inequality does not see \emph{collapse} into the single inequality $u_r \ge 0$;
\item each remaining inequality has one term in $u_r$ and one term per structural column $j$ with $(V_b)_{jp} \neq 0$.
\end{itemize}
The package creates, for each inequality $r$ that is not deterministic, the variables $u_r \ge 0$ and $S_{rj}$, $j \in \Jr_r$,
with the equalities $S_{rj} = \sigma_r(b_{rj} - (a_r^\T X)_j)$ and $u_r = \sum_{j\in\Jr_r} \mub_j S_{rj}$,
and the inequalities \eqref{eq:dual-cone-row} that have a slack term (\code{\_add\_dual\_cone!}).
Restricting $S_r$ to $\Jr_r$ is not a restriction: the other entries are zero for every $X \in \Rules$, as noted after \eqref{eq:slack-rule}.
This is the masking of $S$ by the history that \citet[(4.6)]{KWG2011} write as $S_tP_t$.

\begin{proposition}[Structured dual]
\label{prop:structured-dual}
The structured problem above and the dense problem \eqref{eq:DLDR} have the same feasible rules $X$ and the same objective.
\end{proposition}

\begin{proof}
By \eqref{eq:dual-identity}, the inequalities $S_r\Mom\Wp^\T \ge 0$ are $u_r \ge 0$ and the inequalities \eqref{eq:dual-cone-row} for all rows of all blocks, with $u_r = S_r\mub$.
Those that are dropped are implied by $u_r \ge 0$ and $W_{bp}\mu_b - h_{bp} \ge 0$.
The variables $u_r$ and $S_{rj}$ are determined by $X$ through the equalities.
\end{proof}

The previous version's spike validated this identity numerically on the hydrothermal problem
(objectives equal to a relative difference of $10^{-12}$ on 12 instances),
and the test suite compares both formulations with dense oracles (\file{test/oracles/dense.jl}).
The auxiliary $u_r$ keeps each inequality \eqref{eq:dual-cone-row} as sparse as the covariance allows:
substituting it out would add the $\abs{\Jr_r}$ terms $(\Wp\mub)_p\mub_j S_{rj}$ to every inequality of the row.
Whether to keep it is a question of solver behaviour more than of size, discussed in \cref{sec:computation}.

\subsection{Reporting the bound}

The rule $X$ of \eqref{eq:DLDR} is not a policy: it need not be feasible for any outcome.
The package reports only its value, as the objective bound of the model, and the relative gap between it and the primal value.
With integer decisions, the inner problem is a mixed-integer program, and the bound is the inner solver's bound, which needs no incumbent;
for a linear program, it is the optimal value, or the dual objective value of the inner solver if the solve stopped early.
Since \eqref{eq:DLDR} is a relaxation, its infeasibility proves the infeasibility of \eqref{eq:SP} (under the conditions of \cref{thm:dual-bound}),
while an infeasible primal approximation only shows that no linear rule is feasible:
the package reports \code{INFEASIBLE} only on the evidence of the dual, and \code{LOCALLY\_INFEASIBLE} otherwise.
The bound is \emph{certified} when the moments are exact, $\Mom \succ 0$, and every law spanning several stages is known to satisfy condition~\ref{it:lce} (\cref{rem:lce});
otherwise the result says why it is not (\cref{sec:rejection}).

\section{Sampled decision rules}
\label{sec:sampled}

To be written.

\section{Liftings and piecewise affine rules}
\label{sec:liftings}

To be written.

\section{Uncertainty sets and laws}
\label{sec:laws}

To be written.

\section{Computational notes}
\label{sec:computation}

To be written.

\section{Outlook}
\label{sec:outlook}

To be written.


\appendix
\section{The block robust counterpart}
\label{app:primal}

To be written.

\section{The dual approximation}
\label{app:dual}

To be written.

\section{Moments}
\label{app:moments}

\subsection{Lifted moments}

\begin{proof}[Proof of \cref{prop:lifted-moments}]
With $v_0 = \underline\eta$, the shifted coordinates are $\upsilon_1 = \leta_1 - \underline\eta = \min(\eta - v_0, \Delta_1)$ (on the support, $\eta \ge v_0$)
and $\upsilon_i = \min(\max(\eta - v_{i-1}, 0), \Delta_i)$ for $i \ge 2$, the last one without the upper clamp, which the support makes redundant.
On the piece $i$, that is when $v_{i-1} < \eta \le v_i$,
\[
\upsilon_j = \Delta_j \ (j < i), \qquad \upsilon_i = \eta - v_{i-1}, \qquad \upsilon_j = 0 \ (j > i).
\]
(The atom $\eta = v_0$ of a discrete law, which belongs to no piece, gives $\upsilon = 0$ and contributes nothing below.)
Hence, summing over the pieces,
$\E[\upsilon_i] = \E[(\eta - v_{i-1})\one_i] + \Delta_i\sum_{j > i}\E[\one_j] = m_i + \Delta_i\theta_i$
and $\E[\upsilon_i^2] = s_i + \Delta_i^2\theta_i$.
For $i < j$, $\upsilon_j > 0$ only on the pieces $j' \ge j > i$, where $\upsilon_i = \Delta_i$,
so $\upsilon_i\upsilon_j = \Delta_i\upsilon_j$ everywhere and $\E[\upsilon_i\upsilon_j] = \Delta_i\E[\upsilon_j]$.
Finally $\leta = \upsilon + \underline\eta e_1$ shifts the mean and leaves the covariance unchanged.
\end{proof}

For the uniform law on $[\underline\eta, \overline\eta]$, with $w = 1/(\overline\eta - \underline\eta)$,
the piece moments are $p_i = w\Delta_i$, $m_i = w\Delta_i^2/2$, $s_i = w\Delta_i^3/3$,
and $\theta_i = w(\overline\eta - v_i)$, so
\[
\E[\upsilon_i] = w\Delta_i\Bigl(\frac{\Delta_i}2 + \overline\eta - v_i\Bigr),
\qquad
\E[\upsilon_i^2] = w\Delta_i^2\Bigl(\frac{\Delta_i}3 + \overline\eta - v_i\Bigr),
\qquad
\E[\upsilon_i\upsilon_j] = \Delta_i\,\E[\upsilon_j] \ (i < j).
\]
For instance, for $\mathcal{U}(80, 120)$ with the breakpoint $100$, $\E[\leta] = (80 + 15, 5) = (95, 5)$.
For other continuous laws, $p_i$ comes from the distribution function, and $m_i$, $s_i$ from quadrature of the density on the piece
(\code{\_piece\_moments}, \file{src/lifting/axis_piecewise.jl}).

\subsection{The confidence-ellipsoid normal law}

Let $\eta \sim \mathcal{N}(\mu, \Sigma)$ in $\R^d$ with $\Sigma = LL^\T \succ 0$, and $z = L^{-1}(\eta - \mu) \sim \mathcal{N}(0, I_d)$.
The ellipsoid $\mathcal{E}$ is $\{\norm z^2 \le \rho^2\}$, and $R \defeq \norm z^2 \sim \chi^2_d$, with density
$f_d(r) = r^{d/2-1}e^{-r/2}/(2^{d/2}\Gamma(d/2))$ on $r > 0$.
So $\Prob(\eta\in\mathcal{E}) = \Fchi{d}(\rho^2) = \alpha$ defines $\rho^2 = \Fchi{d}^{-1}(\alpha)$.

\begin{proof}[Proof of \cref{prop:confidence-moments}]
\emph{Mean.} The law of $z$ and the ball are invariant under $z \mapsto -z$, so $\E[z \mid \norm z \le \rho] = 0$ and the mean of $\eta$ is $\mu$.

\emph{Covariance.} The law of $z$ and the ball are invariant under rotations, so $\E[zz^\T \mid \norm z \le \rho] = c\,I_d$ for some $c$,
and taking the trace, $c = \E[R \mid R \le \rho^2]/d$.
Now $r f_d(r) = d\, f_{d+2}(r)$, because $2^{d/2+1}\Gamma(d/2+1) = d\,2^{d/2}\Gamma(d/2)$.
Therefore
\[
\E[R\,\ind{R \le \rho^2}] = \int_0^{\rho^2} r f_d(r)\,dr = d\,\Fchi{d+2}(\rho^2),
\qquad
c = \frac{d\,\Fchi{d+2}(\rho^2)}{d\,\alpha} = \frac{\Fchi{d+2}(\rho^2)}{\alpha} = c_\alpha,
\]
and $\Cov(\eta) = L\,(c_\alpha I_d)L^\T = c_\alpha\Sigma$.
Since $\chi^2_{d+2}$ is stochastically larger than $\chi^2_d$, $\Fchi{d+2}(\rho^2) < \Fchi{d}(\rho^2) = \alpha$, so $c_\alpha < 1$:
truncation shrinks the covariance.

\emph{Conditional expectation.}
Split $\eta = (\eta_h, \eta_f)$ and $\Sigma$ accordingly; let $m(\eta_h) = \mu_f + \Sigma_{fh}\Sigma_{hh}^{-1}(\eta_h - \mu_h)$ and $\Sigma_{f|h} = \Sigma_{ff} - \Sigma_{fh}\Sigma_{hh}^{-1}\Sigma_{hf} \succ 0$.
The Schur complement gives
\[
(\eta - \mu)^\T\Sigma^{-1}(\eta - \mu) = (\eta_h - \mu_h)^\T\Sigma_{hh}^{-1}(\eta_h - \mu_h) + (\eta_f - m(\eta_h))^\T\Sigma_{f|h}^{-1}(\eta_f - m(\eta_h)).
\]
Under the untruncated law, $\eta_f$ given $\eta_h$ is $\mathcal{N}(m(\eta_h), \Sigma_{f|h})$.
Under the truncated law, its conditional density is that density restricted to the ellipsoid
$\{(\eta_f - m)^\T\Sigma_{f|h}^{-1}(\eta_f - m) \le \rho^2 - (\eta_h - \mu_h)^\T\Sigma_{hh}^{-1}(\eta_h - \mu_h)\}$ and renormalized.
Both are symmetric under $\eta_f - m \mapsto -(\eta_f - m)$, so $\E[\eta_f \mid \eta_h] = m(\eta_h)$, which is affine in $\eta_h$.
\end{proof}

The formulas of the previous version, written with incomplete gamma functions, are the same:
$\Fchi{p}(x) = \gamma(p/2, x/2)/\Gamma(p/2)$, so $\rho^2$ solves $\gamma(d/2, \rho^2/2) = \alpha\Gamma(d/2)$
and $c_\alpha = \gamma(d/2+1, \rho^2/2)/(\alpha\Gamma(d/2+1))$.

\section{Size counts}
\label{app:sizes}

To be written.


\bibliographystyle{abbrvnat}
\bibliography{references}

\end{document}
